\chapter{Design and Architecture}
\label{ch:design}

\section{Introduction}
This chapter presents the design of the \projname. It begins with the principles
that governed every subsequent decision, then presents the system architecture
and the alternatives that were considered and rejected. It decomposes the system
into modules, specifies the data design, reports the exploratory analysis that
informed the modelling choices, and then develops in turn the four design
decisions that give the system its character: the feature representation and its
leakage discipline, the horizon-as-feature formulation, the construction of
prediction intervals, and the promotion gate. It closes with the communication
design, the error-handling strategy and the deployment view.

\section{Design Principles}
Six principles were adopted at the outset and applied consistently.

\textbf{P1 - Durable state, disposable compute.} Every machine that runs this
system is destroyed minutes after it starts. The design therefore treats compute
as worthless and state as precious: nothing of value may remain on a runner when
a job ends. This principle is what forces a genuine model registry rather than a
saved file on disk, and it is the reason the architecture can be free.

\textbf{P2 - Decoupling through shared state, not through calls.} No pipeline
invokes another. Their only contract is the feature store and the model registry.
The consequence is that any pipeline can fail, be re-run, be rescheduled or be
rewritten without the others being affected, and that a partial system is still a
working system: if training fails tonight, yesterday's champion continues to
serve.

\textbf{P3 - The past may never see the future.} Every feature attached to time
$t$ must contain only information that existed at $t$; every train-validation
split must be chronological; every rolling statistic must be shifted so that it
ends before the current observation; and no lag may cross a city boundary. This
principle is enforced by construction and verified by unit tests, because target
leakage produces a model that scores beautifully and then fails in production,
and the failure is not visible in the metrics that would normally be trusted.

\textbf{P4 - Nothing ships without proving it is better.} An automated pipeline
that retrains nightly will otherwise deploy whatever it happened to produce.
Promotion is therefore gated on a measured improvement over the incumbent, and
every decision - including every refusal - is recorded.

\textbf{P5 - Auxiliary capability must never break the primary one.}
Explanation, monitoring and forecast logging are valuable but subordinate. Each
is wrapped so that its failure is logged and the forecast is still published. A
system that refuses to forecast because it could not explain itself is worse than
one that forecasts without an explanation.

\textbf{P6 - Publish the working.} Accuracy per horizon, the backtest, the
promotion history and the current drift status are part of the product, not
internal diagnostics. They appear in the same interface as the forecast, because
a public-health signal that cannot be checked will not be trusted.

\section{System Architecture}
The architecture, shown in Figure~\ref{fig:arch}, is an instance of the
Feature-Training-Inference pattern~\cite{ref:fti}. It divides into three
planes.

\fig{architecture}{1.0}{The system architecture: a source plane, a disposable
compute plane, a durable state plane, and the serving layer beneath
them.}{fig:arch}

The \textbf{source plane} is a single free, keyless provider~\cite{ref:openmeteo}
supplying three endpoints: an air-quality endpoint that serves both history and
forecast, a weather forecast endpoint, and a weather archive endpoint derived
from the ERA5 reanalysis~\cite{ref:era5}. The choice of a keyless provider was
not incidental. Because there is no credential to manage, the scheduled workflows
need no secret for ingestion, which removes an entire class of operational
failure.

The \textbf{compute plane} consists of scheduled, ephemeral jobs: the hourly
feature pipeline followed by hourly batch inference and forecast publication, the
daily training pipeline with evaluation and monitoring steps, and a manually triggered
backfill. None of them holds state between runs.

The \textbf{state plane} is a managed feature store and model registry. The
feature store holds one feature group keyed on city and timestamp; the registry
holds versioned model bundles tagged with their metrics. Together they are the
sole interface between the pipelines.

The \textbf{serving layer} consists of a container-hosted API, a statically hosted
browser application, and a tool server that exposes the same data over a standard
protocol. Crucially, the serving layer never invokes a pipeline: it reads
documents that a pipeline has already produced.

\subsection{Alternatives Considered}
Three plausible alternative architectures were considered and rejected.

\textbf{A monolithic scheduled script.} A single job that fetches, trains and
predicts would be simpler to write and would need no feature store. It was
rejected because it couples the schedules - ingestion must be hourly and
training need only be daily - because a failure anywhere loses everything, and
because it provides no mechanism for the trained model to outlive the run. It
also reintroduces training-serving skew~\cite{ref:techdebt}, since the features
used for prediction would be recomputed rather than read.

\textbf{A live prediction service.} The model could be loaded into the API and
invoked per request. This would allow any city and any horizon on demand. It was
rejected as the primary path because it makes every page view depend on model
inference, requires an always-on process sized for peak load, and reintroduces
the possibility that a user-facing request fails because a model failed. A
reduced form of it survives, however: the API exposes an on-demand endpoint for
arbitrary coordinates and for what-if simulation, so the capability exists where
it genuinely adds value without the whole product depending on it.

\textbf{One model per city, or one model per horizon.} Both were rejected for the
same reason - artefact proliferation. Twenty-two cities times ten horizons is 220
models to train, version, register, explain and monitor, and each would be
trained on a small slice of the available data. The global horizon-as-feature
design developed in Section~\ref{sec:design-mh} achieves the same coverage with
one artefact.

\section{Module Decomposition}
Figure~\ref{fig:component} shows the decomposition of the Python package into
layers. Dependencies flow strictly downwards.

\fig{component}{1.0}{Module decomposition. Arrows point from a module to the
modules it depends on; no lower layer imports a higher one.}{fig:component}

At the bottom sits the \textbf{data and configuration layer}. A single
configuration module is the sole source of truth for the supported cities, the
filesystem layout, the feature-store object names and the forecast horizon; every
other module imports from it, so adding a city is a one-line change. The
ingestion module wraps the three provider endpoints behind a retrying session.
The index module implements the EPA breakpoints, categories and advisory text.
The store module is a facade over two interchangeable backends, which is what
satisfies NFR-10.

Above it, the \textbf{feature layer} contains exactly two modules: one that turns
raw observations into the engineered feature table, and one that turns the
feature table into the multi-horizon supervised set and splits it chronologically.
Keeping these separate matters, because the first is used by both the feature
pipeline and the inference pipeline, while the second is used only by training
and evaluation.

The \textbf{model layer} contains the registry facade, the leaderboard and
promotion gate, the evaluation suite, the explanation module, and the
sequence-model study, which is deliberately isolated so that its heavy dependency
never reaches the deployment image.

The \textbf{observability layer} holds drift and error monitoring and the
alerting rules, and the \textbf{serving layer} holds the API, the tool
definitions, the advisor, the tool-protocol server and the what-if simulator. The
tool definitions are shared: the advisor and the protocol server call the same
four functions, so a question answered in the dashboard and the same question
asked from an external client cannot disagree.

\section{Data Design}
\label{sec:design-data}
Figure~\ref{fig:datadesign} specifies the data design.

\fig{data-design}{1.0}{Data design: the feature group and its keys, the two
interchangeable storage backends, the published documents, and the structure
of the forecast payload.}{fig:datadesign}

\subsection{The Feature Group}
All engineered features live in one feature group. Its \textbf{primary key} is
the pair (city, timestamp), where timestamp is the observation hour expressed in
epoch seconds; its \textbf{event time} is the corresponding datetime column. This
choice is what makes ingestion idempotent: the hourly pipeline deliberately
fetches the last seven days rather than the last hour, so that lag and rolling
features for the newest observations have sufficient history, and the resulting
overlap is absorbed by the upsert rather than producing duplicates. A failed run
can therefore simply be repeated.

Two configuration decisions were forced by the free tier and are documented here
because they are design decisions rather than accidents. The table format is set
to a time-travel format that materialises server-side, because the platform's
default format requires a client library that the base installation does not
ship. Statistics computation is disabled, because the profiler over a 74-column
table exhausted the memory of the free-tier executor and failed the
materialisation job; statistics are a convenience for the platform's own
interface and are not read by any pipeline, so disabling them costs nothing and
makes the job finish reliably.

\subsection{The Storage Facade}
Every pipeline addresses storage through one module exposing four operations:
save features, read features, list which cities are present, and report the
earliest date held per city. Two backends implement it. The managed feature store
is selected whenever an API key is present; a local columnar store is used
otherwise. Both hold the identical schema and honour the same primary key. This
is a straightforward application of dependency inversion, and it earns its keep
twice: it allows the entire pipeline to be developed and tested on a platform
where the managed client cannot be installed, and it means the choice of vendor
is confined to a single module.

\subsection{Published Documents}
The pipelines publish forecast, evaluation, monitoring, leaderboard, statistics
and global-SHAP documents. Forecasts are committed hourly; evaluation, monitoring
and leaderboard artefacts are refreshed by the daily workflow. These documents are
the only things the serving layer reads: the forecast payload, the evaluation
report, the monitoring report, the leaderboard and the analytics snapshots. Publishing them as
committed files rather than as a live query is what allows the deployed backend
to omit the feature-store dependency entirely, and what guarantees that a page
view never waits for a model.

The structure of the forecast payload is given in the lower half of
Figure~\ref{fig:datadesign}. Each city carries its province and coordinates for
the map, its current observed index and category, 72 hourly points each with a
predicted value, interval bounds and a category, a three-day aggregation, an
alert object, and an explanation object.

\section{Exploratory Data Analysis}
\label{sec:design-eda}
Before any model was designed, the assembled dataset - 696,960 hourly rows
across 22 cities from 1~January~2023 to August~2026 - was analysed. Five
findings shaped the design.

\subsection{The Cities Differ Enormously}
Figure~\ref{fig:eda-city} ranks the cities by mean index over the whole record.
Faisalabad is the most polluted at a mean of 156.9, marginally ahead of Lahore at
151.5; Gilgit is the cleanest at 75.7. The spread is more than a factor of two.

\fig{eda-city-ranking}{0.86}{Mean Air Quality Index by city over the full
record.}{fig:eda-city}

Two design consequences follow. First, five cities have a \emph{mean} above 140,
which is to say that the typical condition in those cities is at or near the
threshold the EPA calls unhealthy; a forecasting service for them is not a
curiosity. Second, the fact that Faisalabad outranks Lahore is a direct argument
for the multi-city scope of the project, since the city that receives essentially
all of the public and academic attention is not in fact the worst affected.

\subsection{Seasonality Is the Dominant Signal}
Figure~\ref{fig:eda-season} shows the monthly profile. The winter peak is
unmistakable: January is the worst month at a mean of 146.8, against 91.5 in
April, the cleanest - a 60\% difference between the extremes of the year, and a
winter-to-summer ratio of about 1.25. This is consistent with the meteorological
mechanism of the smog season: temperature inversions that trap emissions under a
shallow mixing layer.

\fig{eda-seasonality}{0.86}{Monthly seasonality of the Air Quality Index. The
winter smog season dominates the annual cycle.}{fig:eda-season}

This finding produced two design decisions. It motivated the cyclical encoding of
month, so that December and January are adjacent in feature space rather than
maximally distant. And it established, before any model existed, that the system
would experience severe recurring input drift twice a year - which is why drift
monitoring was specified as a requirement and why its output is interpreted with
the seasonal explanation in mind rather than treated as a fault.

\subsection{There Is a Strong Diurnal Cycle}
Figure~\ref{fig:eda-diurnal} shows the average profile across the hours of the
day. The maximum falls in the evening, around 18:00 at a mean of 123.1, and the
minimum in the morning, around 09:00 at 108.1 - a within-day swing of 15 index
points, consistent with boundary-layer behaviour and the evening traffic peak.

\fig{eda-diurnal}{0.86}{Average diurnal profile of the Air Quality
Index.}{fig:eda-diurnal}

This motivated the cyclical encoding of hour, and it also motivated an interface
decision: because the within-day variation is large, an hourly view is genuinely
useful to a user, not merely decorative, since it lets them identify a safer part
of the day rather than writing off the day entirely.

\subsection{Weather Is Informative but Not Individually Decisive}
Figure~\ref{fig:eda-weather} shows the correlation of each weather variable with
the index. The strongest single relationship is with surface pressure, at a
correlation of about 0.32 - meaningful, but far from sufficient on its own.

\fig{eda-weather-correlation}{0.86}{Correlation between weather variables and the
Air Quality Index.}{fig:eda-weather}

This is an argument for a non-linear model with interactions rather than a linear
one. No single weather variable determines air quality; what matters is the
combination - low wind together with high pressure and a cold night is the
inversion signature - and that is precisely the kind of structure a tree ensemble
captures and a linear model does not. The eventual gap between the Ridge and
XGBoost candidates, reported in Chapter~\ref{ch:testing}, confirms this reading.

\subsection{The Target Distribution Is Skewed}
Figure~\ref{fig:eda-dist} shows the distribution of the index over the whole
record: a long right tail, with the bulk of the mass in the moderate-to-unhealthy
range and occasional excursions into the hazardous band.

\fig{eda-aqi-distribution}{0.8}{Distribution of the Air Quality Index across the
whole record.}{fig:eda-dist}

The skew has an important consequence for the interval design: a symmetric
Gaussian interval would be a poor description of the error, which motivated the
use of empirical residual quantiles rather than a parametric interval.

\section{Feature Design}
Figure~\ref{fig:featureeng} specifies the transformation from raw observations to
model inputs.

\fig{feature-engineering}{1.0}{Feature design: 21 raw columns become a 74-column
engineered table, from which a 27-input multi-horizon supervised set is
assembled. The rule at the foot governs every step.}{fig:featureeng}

\subsection{Feature Families}
Four families of engineered feature are constructed, each for a stated reason.

\emph{Cyclical calendar features.} Hour, month and day of week are each mapped to
a sine and cosine pair. Encoding hour as the integer 23 places it maximally far
from hour 0, which is false; the sine-cosine pair places them adjacent on a
circle, which is true.

\emph{Wind as a vector.} Wind direction is a bearing, so 359 degrees and 1 degree
describe almost the same wind while differing maximally as numbers. Speed and
bearing are therefore decomposed into orthogonal components, which the model can
reason about linearly.

\emph{Lagged values.} The index and both particulate species are lagged by 1, 3,
6, 12 and 24 hours, since pollutant series are strongly autocorrelated and the
24-hour lag additionally captures the diurnal cycle.

\emph{Rolling statistics.} Mean, standard deviation and maximum over 6-hour and
24-hour windows summarise recent level, volatility and peak.

\subsection{The Leakage Discipline}
Every rolling window is shifted by one observation before it is computed, so that
it ends at the previous hour and can never include the current one. Every lag and
rolling statistic is computed within a single city, by grouping before
engineering, so that one city's history can never leak into another's. Both
properties are asserted by unit tests: one checks that each city independently
shows exactly 24 missing values in its 24-hour lag column, which can only be true
if the groups were processed separately, and another checks that the anchor value
at horizon $h$ equals the observed value $h$ steps earlier.

\subsection{The Target}
The target is the EPA index. Where the provider supplies its own correctly
averaged index value, that is used directly; where it is missing, the index is
computed from first principles from particulate matter. Gaseous pollutants are
deliberately excluded from the computed fallback, because the EPA defines their
breakpoints in volumetric mixing ratios while the provider reports mass
concentrations, and applying one to the other without a molar-mass conversion
inflates the sub-index absurdly. Since the index is the maximum over sub-indices,
a single wrong sub-index would corrupt every affected row. The exclusion is
tested explicitly.

\section{Multi-Horizon Forecasting Design}
\label{sec:design-mh}
The central modelling decision is illustrated in Figure~\ref{fig:multihorizon}.

\fig{multihorizon}{1.0}{Direct multi-horizon forecasting with the horizon as an
input feature: one global model serves 22 cities and ten lead
times.}{fig:multihorizon}

\subsection{The Two Kinds of Input}
A forecast issued at time $t$ for the hour $\tau = t + h$ can draw on two
distinct kinds of information, and the design separates them explicitly.

\emph{Target-time features} describe the hour being predicted and are known in
advance: the weather at $\tau$, which is available as a forecast; the calendar at
$\tau$, which is deterministic; and the location, which is fixed. There are 20 of
these.

\emph{Anchor-state features} describe the observed pollution at the moment the
forecast is issued: the current index, the 6-hour and 24-hour index means, the
recent index volatility, and the current particulate concentrations. There are
six of these, and they are obtained by shifting each series by $h$, so that the
row describing $\tau$ carries the state as it was $h$ hours earlier.

The horizon $h$ itself is the twenty-seventh input. The model therefore learns
\[
\mathrm{AQI}(\tau) \;=\; f\big(\text{target-time features at } \tau,\;
\text{anchor state at } \tau - h,\; h\big).
\]

\subsection{Why the Horizon Is a Feature}
Making $h$ an input rather than training one model per horizon has four
consequences, all favourable. There is one artefact to version, register, explain
and monitor, rather than ten or 220. Every horizon is trained on the whole
dataset rather than on a tenth of it, so the long horizons - which are the
sparsest and the hardest - borrow strength from the short ones. The model learns
how uncertainty grows with distance as a smooth function of $h$ rather than
memorising ten unrelated mappings, so an unmodelled lead time can be served by
interpolation. And unlike a recursive strategy~\cite{ref:multistep}, the model is
never fed its own output, so errors do not compound across steps.

The design also incurs an honest cost, which Chapter~\ref{ch:testing} quantifies:
at training time the model sees the weather that actually occurred at $\tau$,
whereas at inference it sees the forecast weather for $\tau$. The validation
score is therefore a mild over-estimate of operational skill.

\subsection{The Supervised Set and Its Split}
Ten horizons - 1, 2, 3, 6, 12, 24, 36, 48, 60 and 72 hours - are materialised.
Each observation therefore contributes up to ten supervised rows differing only
in $h$ and in the target. The resulting set is subsampled to 200,000 rows with a
fixed seed, which keeps a full four-candidate training run inside the scheduler's
time limit while retaining ample data.

The split is strictly chronological: the most recent fifth of the time range
becomes validation, and everything before it becomes training. Random
partitioning is not used at any point, in accordance with the
literature~\cite{ref:bergmeir} and with principle P3.

\section{Prediction Interval Design}
\label{sec:design-intervals}
A point forecast alone would violate NFR-4. Intervals are constructed in the
spirit of split conformal prediction~\cite{ref:inductive,ref:leiconformal}: the
model is fitted on the training window, its residuals are computed on the
disjoint validation window, and the interval is formed from the empirical
quantiles of those residuals.

Two refinements are applied. The quantiles are computed \emph{per horizon}, so
that the interval widens with lead time exactly as the model's real errors do; a
single global interval would be far too wide at one hour and far too narrow at 72.
And the quantiles are empirical rather than parametric, which respects the skew
observed in Section~\ref{sec:design-eda} instead of assuming a symmetric error
distribution. The tenth and ninetieth percentiles are used, giving a nominal 80\%
interval. At inference the bounds are clipped to the valid index range, and a
horizon that was not explicitly modelled takes the quantiles of the nearest
modelled horizon.

The construction's assumption - exchangeability between the calibration and
deployment periods - is not strictly satisfied by a seasonal, non-stationary
series, so the 80\% figure is a well-motivated estimate rather than a guarantee.
Chapter~\ref{ch:testing} states this limitation explicitly rather than claiming
coverage the construction does not earn.

\section{Model Selection and Promotion Design}
\label{sec:design-gate}
Four candidate families are scored in every run: a persistence baseline, which
predicts that the index at $t+h$ equals the index at $t$; a standardised Ridge
regression; a random forest; and a gradient-boosted ensemble. The baseline is
scored but is never eligible for promotion; its role is to establish that the
learned models have skill at all~\cite{ref:fpp}. The linear model's role is
diagnostic, since the gap between it and the ensembles measures how much of the
relationship is non-linear.

The best fitted candidate by validation RMSE becomes the challenger, and is
promoted only if it improves on the reigning champion's validation RMSE. The
comparison margin is zero, which means an exact tie is a refusal: a new model
never ships merely because it is newer. Every run appends an entry to a
persistent leaderboard recording the version, the timestamp, the winning family,
its metrics, the metrics of every candidate considered, the previous champion's
score and the decision. The design and its observed behaviour are developed in
Section~\ref{sec:impl-gate}.

\section{Monitoring Design}
The monitoring design pairs two measurements that are individually insufficient.

\emph{Input drift} is measured by the Population Stability Index over seven
features - the index itself, both particulate species, temperature, humidity,
wind speed and surface pressure - comparing the most recent fourteen days against
everything older. Deciles of the reference distribution define the bins, and the
conventional thresholds of 0.10 and 0.25 classify each feature as stable,
moderate or significant. PSI is non-parametric, cheap, and computed per feature so
that the responsible variable is named rather than merely alarmed about.

\emph{Realised error} is measured directly. Every forecast point issued is
appended to a log with the hour it was issued for; once the feature pipeline has
ingested the observation for that hour, the log is joined to it on city and target
hour, and MAE, RMSE and the five largest misses are computed over the joined rows.

The pairing is the point. Input drift alone cannot distinguish a harmless
seasonal transition from a genuine degradation, and for a system whose dominant
signal is an annual cycle, a drift alarm is the expected condition for several
months of the year. Realised error alone is a lagging indicator: it can only
score a forecast once the future it described has arrived. Read together, they
separate the two cases - drift with stable error is seasonality that the nightly
retrain has absorbed; drift accompanied by rising error is decay.

\section{Communication Design}
The system exposes one HTTP interface, summarised in Table~\ref{tab:api}. Its
design follows a single rule: endpoints that read a published document are cheap
and always available, while endpoints that invoke the model are separate,
explicitly labelled, and permitted to fail without affecting the others.

\begin{xltabular}{\textwidth}{P{4.5cm} Q{1.6cm} Y}
\caption{The HTTP interface.}\label{tab:api}\\
\toprule
\bfseries Endpoint & \bfseries Kind & \bfseries Purpose \\
\midrule
\endfirsthead
\toprule
\bfseries Endpoint & \bfseries Kind & \bfseries Purpose \\
\midrule
\endhead
\bottomrule
\endlastfoot
/api/health & document & Liveness, and whether a forecast document is present. \\
/api/categories & static & The six index categories, their bounds and their
advisory text, for the legend. \\
/api/cities & static & The supported cities with province and coordinates. \\
/api/predictions & document & The complete forecast payload for every city. \\
/api/predictions/\{city\} & document & One city's forecast. \\
/api/leaderboard & document & The promotion history and the current champion. \\
/api/evaluation & document & Per-horizon metrics and the walk-forward backtest. \\
/api/monitoring & document & Drift status and realised forecast error. \\
/api/explain/\{city\} & document & The stored attribution and its sentence. \\
/api/predict & model & An on-demand forecast for an arbitrary coordinate pair. \\
/api/whatif/defaults & model & Current real values for the scenario sliders. \\
/api/whatif & model & Baseline versus scenario under overridden drivers. \\
/api/chat & external & The grounded language-model advisor. \\
\end{xltabular}

The frontend is a single-page application that fetches these endpoints over
HTTPS. Because the two are deployed to different origins, the API permits
cross-origin requests; this is safe here because every endpoint is a read of
public, non-sensitive data. When no backend address is configured at build time,
the same frontend code reads committed snapshots shipped alongside the bundle and
hides the two capabilities that genuinely require a live model, which is the
graceful degradation required by NFR-11.

Figure~\ref{fig:sequence} shows how these pieces interact over time, and makes
the essential asymmetry visible: the left-hand portion is a scheduled batch job
that may take minutes, while the right-hand portion is a file read that takes
milliseconds, and nothing connects them except a document.

\fig{sequence-inference}{1.0}{Sequence of a scheduled forecast refresh, followed
by an independent dashboard request.}{fig:sequence}

\section{Error Handling and Recovery Design}
Failure is treated as ordinary rather than exceptional, and is handled
differently at each layer.

\emph{At the network boundary}, transient provider failures are retried
automatically with exponential backoff on the status codes that indicate
throttling or a temporary server fault. Without this, a single momentary error
would fail an unattended hourly job.

\emph{Within a multi-city loop}, a failure is contained. If one city's ingestion
or forecast raises, the error is logged and the loop continues, so that
twenty-one cities are still served rather than none. The backfill goes further
and writes each city as it completes, so an interruption loses at most one city's
work.

\emph{At the storage boundary}, writes are idempotent upserts on the primary key
and are submitted without blocking on server-side materialisation. This design
came directly from a real failure: a blocking write waited indefinitely on a
stalled free-tier job and froze an entire backfill run. Because writes are
idempotent and the resume logic is history-aware, simply re-running the job is
always a safe recovery.

\emph{At the model boundary}, the promotion gate is itself an error-handling
mechanism. If a training run produces a bad model - because of a corrupt window,
an ingestion gap or any other cause - the gate refuses it and the previous
champion continues to serve. There is no path by which an unattended run can
degrade the service.

\emph{At the auxiliary boundary}, each subordinate capability is wrapped so that
its failure is logged and execution continues, in accordance with principle P5.
The monitoring endpoint carries this furthest, with three levels of fallback: it
serves the committed snapshot if present, computes the report live if not, and
returns a clear placeholder if neither is possible, so that the tab reports its
own state rather than showing an error.

\section{Deployment View}
Figure~\ref{fig:deployment} shows the deployment topology and the two modes in
which the system can serve.

\fig{deployment}{1.0}{Deployment view: the repository, the scheduled runners, the
container-hosted API, the statically hosted frontend, and the two serving
modes.}{fig:deployment}

Four nodes participate. The \textbf{repository} holds the source and, uniquely,
the published documents and the model bundle, which the hourly and daily workflows commit
back to it. The \textbf{runners} execute the pipelines and are the only nodes that
ever contact the feature store. The \textbf{container host} runs the API from a
deliberately slim image that omits the feature-store and deep-learning
dependencies entirely, because it serves committed documents and a bundled model
rather than querying live state. The \textbf{static host} serves the compiled
frontend from the edge with no server process at all.

The two serving modes are a direct expression of NFR-11. In \emph{dynamic} mode,
which is how the system is currently deployed, the frontend is built with the
backend's address and fetches everything live, so the full capability set -
including what-if simulation and the advisor - is available. In \emph{static}
mode, no backend address is configured, the same frontend code reads the
committed snapshots, and the two capabilities that require a live model are
hidden rather than shown broken. The second mode is not a degraded accident but a
designed fallback: it guarantees that the forecast remains publicly readable even
if the container host is unavailable.

\section{Summary of Design and Architecture}
This chapter set out the six principles that governed the design - durable state
with disposable compute, decoupling through shared state, strict prevention of
temporal leakage, gated promotion, subordination of auxiliary capability, and
publication of the system's own working. It presented the three-plane
Feature-Training-Inference architecture and the alternatives rejected in
reaching it, decomposed the package into layers, and specified the data design
around a single upsert-keyed feature group and published forecast, evaluation,
monitoring, statistics and SHAP documents. It
reported the exploratory analysis that motivated the cyclical encodings, the
non-linear model class, the empirical interval construction and the drift
monitoring. It then developed the four decisions that define the system: the
leakage-controlled feature representation, the horizon-as-feature formulation,
per-horizon empirical intervals, and the promotion gate. Finally it specified the
HTTP interface, the analytics snapshot fallback, the layered error-handling strategy and the deployment topology
with its two serving modes. The next chapter describes how this design was
implemented.
